\section{Constraint-Aware MLP}
Consider a semilinear parabolic PDE
\begin{equation}
\partial_tu+\mathcal Lu+f(t,x,u,z)=0,\qquad u(T,x)=g(x),
\label{eq:pde}
\end{equation}
with gradient-like state $z=\sigma(t,x)^\top\nabla_xu$. MLP recursively estimates $Y=(u,z)$ at random space--time points and reuses those estimates inside $f$. The relevant failure mode is not merely that $\widehat Y$ has large variance; it is that the nonlinear map can destroy cancellation. For example, if the exact gradient obeys $\norm{z}\le 4$ but a finite-sample estimate has norm $12$, then the quadratic HJB driver $f(z)=-\tfrac12\norm z^2$ turns a factor-three state error into a factor-nine driver magnitude. Once this value enters a higher Picard level, a final-output correction is too late.

\begin{figure}[t]
\centering
\setlength{\fboxsep}{4pt}
\fbox{\parbox{0.92\columnwidth}{\centering
raw child $\widehat Y$ $\;\longrightarrow\;$ certified retraction $\mathcal R$ $\;\longrightarrow\;$ nonlinear $f(t,x,\mathcal R(\widehat Y))$ $\;\longrightarrow\;$ parent state}}
\caption{The correction is inserted at the nonlinear reuse point, not only at the final output.}
\label{fig:interface}
\end{figure}

\subsection{Two recursive retractions}
Let $Y_1,\ldots,Y_B$ be sibling child-state estimates and suppose PDE theory certifies $Y^\star(t,x)\in\mathcal C(t,x)$.

\paragraph{Samplewise IR.}
For nonempty closed convex $\mathcal C$,
\begin{equation}
\widetilde Y_i=\Pi_{\mathcal C}Y_i.
\label{eq:samplewise}
\end{equation}
This is local, introduces no cross-sample coupling, and makes the smallest Euclidean correction to each child.

\paragraph{Batch-IR.}
For a gauge set $\mathcal C=\{y:\gamma_{\mathcal C}(y)\le1\}$, define
\begin{equation}
\alpha_B=\min\!\left\{1,\frac{1}{\max_i\gamma_{\mathcal C}(Y_i)}\right\},
\qquad \widetilde Y_i=\alpha_BY_i.
\label{eq:batch}
\end{equation}
For a ball, $\alpha_B=\min\{1,R/\max_i\norm{Y_i}\}$. The rule preserves relative directions and ratios in the batch, but it couples siblings and generally changes their empirical mean.

\begin{table}[t]
\caption{Same feasibility principle, different statistical geometry.}
\label{tab:methods}
\centering\small
\begin{tabular}{@{}lcc@{}}
\toprule
& Samplewise & Batchwise\\
\midrule
map & $\Pi_{\mathcal C}Y_i$ & $\alpha_BY_i$\\
coupling & none & shared\\
nearest point & yes & no\\
theory & projection & structural rescue\\
risk & proj. bias & coupling/bias\\
\bottomrule
\end{tabular}
\end{table}

\subsection{What can be guaranteed without a new MLP theory?}
The key deterministic geometry is the projection inequality: for every feasible $Y^\star\in\mathcal C$,
\begin{equation}
\norm{\Pi_{\mathcal C}Y-Y^\star}^2
\le \norm{Y-Y^\star}^2-\dist(Y,\mathcal C)^2.
\label{eq:pyth}
\end{equation}
Thus samplewise projection removes at least the squared overshoot component of child-state error.

\begin{proposition}[Target preservation]
\label{prop:target}
If $Y^\star(t,x)\in\mathcal C(t,x)$ and the retraction is the identity on feasible states, then constraint-aware evaluation leaves the exact generator unchanged:
\[
f\bigl(t,x,\mathcal R(Y^\star(t,x))\bigr)=f(t,x,Y^\star(t,x)).
\]
Hence the correction stabilizes off-target stochastic states rather than changing the intended PDE solution.
\end{proposition}
\begin{proof}
Feasibility gives $\mathcal R(Y^\star)=Y^\star$.
\end{proof}

For samplewise convex projection, this can be strengthened under the standard assumptions used by gradient-dependent MLP theory: defining $\bar f(t,x,y)=f(t,x,\Pi_{\mathcal C(t,x)}y)$ yields an ordinary MLP problem with the same feasible solution, and nonexpansiveness of projection does not increase the state Lipschitz constant. We give the short inheritance argument in Appendix~\ref{app:convergence}; it is not a headline theorem here.

\paragraph{Work order.}
Projection onto the subspaces, balls, and normalized ellipsoids used below costs $O(d)$ per state. Forming a batch gauge and rescaling $B$ $d$-dimensional states costs $O(Bd)$, the same arithmetic order as reading the batch. For Samplewise IR, the inherited convergence result therefore leaves the usual polynomial complexity exponents unchanged under the same standing assumptions. For Batch-IR we claim only unchanged per-update arithmetic order; a sharp cost-to-accuracy theory must account for coupling and is future work.

